\begin{frame}{References}
    \vspace{0.6em}
    These slides have been adapted from
    \begin{enumerate}
        \item Elvis Saravia \& Mausam, \href{https://www.cse.iitd.ac.in/~mausam/courses/col772/spring2020/}{COL772: Natural Language Processing, IITD}
        \item \href{https://www.promptingguide.ai/}{https://www.promptingguide.ai/}
        \item Tatsunori Hashimoto, \href{https://web.stanford.edu/class/cs224n/}{CS224N: Natural Language Processing with Deep Learning, Stanford}
    \end{enumerate}
\end{frame}

\begin{frame}{References}
    \vspace{0.6em}
    \textbf{Key Papers and Resources}
    \begin{itemize}
        \item[\textcolor{red}{$\triangleright$}] Brown et al., 2020. \textit{Language Models are Few-Shot Learners (GPT-3)}
        \item[\textcolor{red}{$\triangleright$}] Wei et al., 2022. \textit{Chain-of-Thought Prompting Elicits Reasoning in LLMs}
        \item[\textcolor{red}{$\triangleright$}] Zhang et al., 2021. \textit{DPR: Dense Passage Retrieval for Open-Domain QA}
        \item[\textcolor{red}{$\triangleright$}] Lewis et al., 2020. \textit{Retrieval-Augmented Generation for Knowledge-Intensive NLP}
        \item[\textcolor{red}{$\triangleright$}] Schick \& Schütze, 2021. \textit{It's Not Just Size That Matters: Small Language Models Are Also Few-Shot Learners}
    \end{itemize}

    \vspace{0.6em}
    \textbf{Online Resources}
    \begin{itemize}
        \item[\textcolor{red}{$\triangleright$}] \href{https://platform.openai.com/docs/guides/prompt-engineering}{OpenAI Prompt Engineering Guide}
        \item[\textcolor{red}{$\triangleright$}] \href{https://python.langchain.com/docs/}{LangChain Documentation}
        \item[\textcolor{red}{$\triangleright$}] \href{https://docs.llamaindex.ai/}{LlamaIndex Documentation}
    \end{itemize}
\end{frame}
